\section{Discussion}
\label{sec:discussion}

\subsection{Interpreting the RQ1/RQ2 contrast}
The central empirical fact of this paper is not either number in
Section~\ref{sec:evaluation} alone, but that they point in opposite
directions under an otherwise identical mechanism: the same refinement
loop, the same proposer model, the same iteration budget, the same
sandboxing discipline, raises acceptance rate on a like-for-like
replication of the original study's objective (RQ1) and simultaneously
underperforms a simple static generator on a differential divergence
objective (RQ2). This is evidence that ``LLM-guided refinement helps''
is not a property of the pipeline in general -- it is a property of the
pipeline \emph{given a specific objective}, and the objective matters more
than we expected going in. Acceptance rate against a single decoder is a
smooth, easy-to-improve-incrementally signal: almost any perturbation that
makes a document more JSON-shaped moves the number in the right direction,
and an LLM has presumably seen enormous amounts of well-formed JSON during
training. Cross-implementation divergence is a much harder target: it
requires a document that is \emph{simultaneously} well-formed enough to
clear one implementation's more lenient path and malformed enough
(in one specific, narrow way) to trip a stricter one, which is a
needle-in-haystack search over a much smaller region of the input space
than ``more JSON-like'' rewards.

We tested, rather than speculated about, the obvious alternative
explanation: that RQ2's gap is mostly sandbox-driven generation overhead
(the LLM cannot import a JSON library and must hand-build JSON text, which
is failure-prone) rather than a genuine targeting shortfall. Directly
permitting \texttt{import json} (Section~\ref{sec:ablation-json}) eliminated
the generation-overhead problem almost entirely -- every ablation run
reached at or near 100\% \texttt{schema\_evaluated} -- yet the divergence
rate did not move: 2.20\% mean, statistically indistinguishable from the
constrained sandbox's 2.93\% ($p \approx 0.69$), both still an order of
magnitude below the static baseline. The RQ2 gap is not generation
overhead at all; it is entirely a targeting problem. This is a stronger and
more precise conclusion than our first attempt at explaining it
(restricting the comparison to schema-evaluated documents only, which could
not fully separate the two effects).

Having localized the problem to targeting, the direct next question is
whether more informative feedback about the targeting problem specifically
closes it. It did not -- and did measurably worse. Feeding back
perturbation categories (Section~\ref{sec:category-feedback}) produced a
0.95\% mean divergence rate, significantly below the score-only prompt's
2.93\% ($p \approx 0.032$, both overall and restricted to schema-evaluated
documents, ruling out proposal-failure noise as the explanation). Taken
together with the ablation, this suggests the limiting factor is not that
the proposer lacks information about \emph{what} tends to work -- it is
given exactly that information twice now, once implicitly (via the score)
and once explicitly (via categories) -- but a harder difficulty actually
implementing an effective divergence-hunting generator in code, one that
additional textual guidance does not straightforwardly fix and may compete
with by adding more for the proposer to reason about at once. We report
this as a genuine negative data point against the intuitive assumption
that richer feedback is strictly better, not as a settled explanation of
why: at $n{=}5$ we can establish the direction of the effect, not its
mechanism.

\subsection{The value of reporting a corrected prompt, not just a final one}
Section~\ref{sec:rq2-prompt}'s prompt fix is itself worth discussing as a
methodological point independent of its numeric effect. The unmodified
prompt drove up an unrelated metric (per-path rejection counts) while
completely failing at the intended one (divergence) across three full
seeded runs -- a proposer optimizing ``move the numbers shown to me''
when the numbers shown do not clearly identify which one is the actual
objective. This mirrors, on a different metric, exactly the concern the
original study's own discussion raised about acceptance-rate
over-optimization crowding out boundary-probing behavior. We treat this as
a reason to report negative or confounded intermediate results explicitly
in future work on prompt-driven refinement objectives, rather than only
the final, working design -- the failure mode is informative on its own,
and silently iterating past it would have understated how easy it is for
an LLM-guided loop to optimize the wrong thing when the prompt does not
make the objective unambiguous.

\subsection{Threats to validity}
\textbf{Construct validity.} RQ2's divergence oracle can only detect
disagreement between the four specific paths instrumented; it says nothing
about whether any individual path is ``correct'' relative to a
specification, only that they disagree. RQ3's ground-truth oracle is
narrower still, scoped to one field (\texttt{id}) chosen because it is the
schema's only bare JSON-number token -- a deliberate, evidenced choice
(Section~\ref{sec:methodology}) after an earlier draft of this design
named a string-typed field that cannot exhibit numeric precision loss at
all.

\textbf{Internal validity.} The $n{=}5$ RQ2 comparison used a corrected
prompt that itself required an intermediate failed attempt to discover
(Section~\ref{sec:rq2-prompt}); we validated the fix on a held-out seed
before committing the full five-seed budget to it specifically to guard
against tuning the prompt to the same data used to report the result, but
we did not attempt further prompt variants beyond the one fix, so we
cannot rule out that a different, equally principled prompt design closes
more of the RQ1/RQ2 gap than the one reported here.

\textbf{Statistical validity.} $n{=}5$ per arm is a small sample, chosen
under a fixed, modest API budget and reported as exactly that throughout,
following the original study's own precedent for its smaller-sample
ablation. Complete rank separation at $n{=}5$ gives an exact $p \approx
0.00794$, which is the smallest attainable exact $p$-value this design can
produce, not a measure of effect size independent of sample size; we
report the underlying per-run rates in Tables~\ref{tab:rq1}, \ref{tab:rq2},
\ref{tab:ablation}, and \ref{tab:category-feedback} for exactly that
reason, and a larger-$n$
replication, budget permitting, is the direct way to strengthen this
result (Section~\ref{sec:future-work}).

\textbf{External validity.} All results are on JSON, on four specific,
pinned-version Dart deserialization paths, with one LLM
(\texttt{gpt-4.1-mini}) at one temperature setting. We make no claim that
the magnitude of the RQ1, RQ2, or follow-up effects, or the RQ3 saturation
pattern's exact threshold behavior, generalizes to other formats, other
Dart packages or package versions, or other proposer models.

\subsection{When this approach is, and is not, the right tool}
The differential divergence oracle (RQ2) is only informative when multiple
independent implementations of the same specification genuinely exist and
are in real concurrent use, which is true of Dart JSON deserialization
specifically because the ecosystem has several competing, actively
maintained codegen frameworks with no single dominant standard. A target
with one canonical implementation and no independent alternative has no
comparison partner for this oracle at all, and the numeric round-trip
oracle (RQ3) or an exception-contract oracle (Section~\ref{sec:evaluation-defects}'s
private-exception classification) would be the applicable choice instead.

\subsection{Future work}
\label{sec:future-work}
Several extensions follow directly from what this study measured but did
not have budget to pursue further. First, since neither the score-only
prompt nor the richer category-feedback prompt
(Section~\ref{sec:category-feedback}) closed the RQ2 gap -- and the latter
did measurably worse -- the open question is no longer "what feedback
helps" but whether the proposer's difficulty is better addressed
structurally: for instance, letting it retain and mutate its
\emph{previous} iteration's generator source rather than rewriting from
scratch each time (this pipeline's loop, reused unmodified from the
original, always discards the prior proposal), or giving it a small library
of worked examples of divergence-causing perturbations applied to an
unrelated schema, rather than more description of the current one. Second,
a larger-$n$ replication of RQ1, RQ2, and both
Section~\ref{sec:ablation-json}/\ref{sec:category-feedback} follow-ups,
budget permitting, would tighten the exact $p$-values reported here past
what $n{=}5$ can produce, the same way the original study's own $n{=}15$
result tightened its earlier smaller-sample finding. Third, the qualitative
union side study (Section~\ref{sec:evaluation-union}) found a real
exception-type difference between \texttt{freezed} and
\texttt{json\_serializable} on thirteen hand-picked documents but was not
run through this paper's automated oracle, generator, or seeded-campaign
machinery at all; building a Payload-shaped generator and harness endpoint
and re-running the RQ2-style baseline-vs-refined comparison on it would
establish whether this exception-type difference (or a genuine value
divergence, not observed in the small hand-picked sample) appears at the
rate and statistical rigor the main Record schema's findings do, rather
than resting on thirteen cases.
